\section{Introduction}
\IEEEPARstart{C}{ooperative} target protection requires a team to respond to several threats at once. A defender approaching one attacker must account for the target behind that attacker, nearby teammates, and competing threats elsewhere. As vehicles are intercepted and threats move between targets, the useful relationships change. A controller trained for a small engagement should reuse these relationships when more defenders, attackers, or protected assets enter the scene. This requirement motivates policies that coordinate through entity relationships rather than a fixed population layout.

Entity-based multi-agent reinforcement learning provides a suitable starting point. Shared set encoders accept variable numbers of entities, while attention extracts information for each agent~\cite{zaheer2017deepsets,lee2019settransformer}. REFIL further encourages reusable interactions through randomized entity-wise value factorization~\cite{iqbal2021refil}. The remaining design question concerns the representation supplied to the controller: how should an individual defender interpret an attacker together with its surrounding teammates and targets, and retain the parts of that context that matter for its next action?

We address this question through \emph{receiver-conditioned, multi-stage relational encoding}. A shared attention block repeatedly updates the entity representations within one physical decision. Each update operates on the context produced by the preceding update. A reader conditioned on the defender's own state accesses every intermediate representation, and a learned fusion combines these reads into the decision feature. Context construction and reading form one encoder: the loop develops relations among entities, while the reader expresses those relations from the perspective of the acting defender. Keeping intermediate reads lets the controller use early context alongside later refinements.

This computation has a different role from memory across physical time. An attacker's target assignment is private, although its position and velocity are observed. The policy therefore retains a recurrent state to interpret motion over successive observations. We use the current self state to query the relational encoder and place temporal recurrence downstream. The resulting architecture, \textbf{Looped Entity Attention Fusion (LEAF)}, augments REFIL's feature with the fused relational read and then updates the defender's GRU once. Network parameters are shared across entities, defenders, and processing rounds. Training samples one to four rounds, and the deployed policy uses four.

Our evaluation follows the intended control problem. In SkyHAD~\cite{skyhad2026}, policies learn with at most ten vehicles per team and three targets, then execute with frozen parameters on 24 configurations reaching 50 vehicles per team and 12 targets. We compare LEAF with REFIL, QMIX-Atten, DCG, SPECTra, ALMA, and a rule-based controller. Reward is the negative of accumulated target damage. LEAF attains the highest group mean reward among these methods in the in-distribution, target-count, team-size, and joint-extrapolation groups. Its gains over REFIL are accompanied by fewer friendly-collision deaths and more attacker interceptions.

The experiments connect this outcome to the representation design. Trained comparisons examine reading before and after contextual processing, placing memory in the reader, and retaining the REFIL host. Execution at the four depths used in training characterizes how refinement affects the frozen policy. Normal forward records then show which intermediate reads supply its decisions, while common-scene trajectories connect interception behavior to target protection. The contributions are:
\begin{enumerate}
\item A receiver-conditioned relational encoder that couples shared iterative entity updates with fusion of intermediate reads, integrated with recurrent cooperative value learning.
\item A population-transfer evaluation covering changes in vehicle and target counts, with learning curves, two operational metrics, and trained architectural comparisons across three training seeds.
\item An empirical account of the controller's operation that connects contextual processing and intermediate-read use to interception and target-protection outcomes.
\end{enumerate}
